\section{No elementary conjugacy in the subcritical window}
\label{a:sec:ht}

This section makes the obstruction precise. The logical skeleton is short:
\begin{enumerate}[label=(\roman*),leftmargin=*]
\item every elementary function satisfies an algebraic differential equation;
\item differential algebraicity survives compositional inversion;
\item by a theorem of Becker and Bergweiler, a solution of Schröder's equation for
      a rational map of degree at least two can be differentially algebraic only
      when the underlying fixed point is \emph{repelling};
\item for $r\in(0,1)$ the fixed point $0$ of $f_r$ is attracting.
\end{enumerate}
Hypertranscendence of $\psi_r$ at every subcritical parameter follows at once ---
with no exceptional cases to exclude, because the attracting window contains no
candidates at all. \Cref{a:sec:scope} then explains what this
does \emph{not} settle.

\subsection{Differential algebraicity, hypertranscendence, elementarity}
\label{a:sec:dadefs}

\begin{definition}[Differentially algebraic; hypertranscendental]
\label{a:def:da}
A function $g$, holomorphic on a domain, or a formal power series, is
\emph{differentially algebraic over $\mathbb C(z)$} if there is a non-zero
polynomial $P$ in $n+2$ variables, for some $n\ge0$, with
\begin{equation}
\label{a:eq:dadef}
P\bigl(z,\,g(z),\,g'(z),\,\dots,\,g^{(n)}(z)\bigr)\equiv0 .
\end{equation}
Otherwise $g$ is \emph{differentially transcendental}, or
\emph{hypertranscendental}.
\end{definition}

\begin{lemma}[Elementary implies differentially algebraic]
\label{a:lem:elemda}
Every elementary function in the sense of Liouville --- any function reachable from
$\mathbb C(z)$ by a finite tower of algebraic operations, exponentials and
logarithms --- is differentially algebraic. Consequently a hypertranscendental
function is not elementary.
\end{lemma}

\begin{proof}[Proof sketch]
Each storey of a Liouvillian tower preserves differential algebraicity: algebraic
extensions do so by elimination; if $g$ is differentially algebraic then so are
$\ee^{g}$ and $\log g$, their defining first-order relations combining with the
equation for $g$; and the class is closed under the field operations and under
composition. These closure properties are classical; see Rubel's survey
\cite{Rubel1989}.
\end{proof}

\begin{lemma}[Closure under compositional inversion]
\label{a:lem:invda}
Let $g$ be holomorphic near $0$ with $g(0)=0$ and $g'(0)\neq0$, and let $g^{-1}$
denote its local compositional inverse. Then $g$ is differentially algebraic over
$\mathbb C(z)$ if and only if $g^{-1}$ is (Moore \cite{Moore1896}; see also
Boshernitzan and Rubel \cite{BoshernitzanRubel1986}).
\end{lemma}

\Cref{a:lem:invda} matters because the literature states Schröder-equation results
in two equivalent but opposite conventions, and conflating them is an easy way to
misquote a theorem. Our $\psi_r$ is the \emph{linearising coordinate}:
$\psi_r(f_r(z)) = r\psi_r(z)$. The function classified by Becker and Bergweiler ---
and by Ritt before them --- is the \emph{linearising parametrisation}
$\varphi_r = \psi_r^{-1}$, which satisfies
\begin{equation}
\label{a:eq:schroderparam}
\varphi_r(r\,t) = f_r\bigl(\varphi_r(t)\bigr),
\qquad
\varphi_r(0)=0,\quad\varphi_r'(0)=1 .
\end{equation}
By \cref{a:lem:invda}, hypertranscendence of either function transfers to the other.

\subsection{The Becker--Bergweiler theorem, stated precisely}
\label{a:sec:BB}

The result needed is the 1995 classification of differentially algebraic Schröder
functions \cite{Becker1995}; we quote it in the form given in Fernandes' survey of
the Schröder, Böttcher and Abel equations \cite[Thm~3.1]{Fernandes2021}.

\begin{theorem}[Becker--Bergweiler \cite{Becker1995}, as restated in
{\cite[Thm~3.1]{Fernandes2021}}]
\label{a:thm:BB}
Let $R$ be a rational map of degree at least two with $R(0)=0$ and multiplier
$s = R'(0)$, where $s\neq0$ and $s$ is not a root of unity; for $|s|=1$ the
Schröder function is understood as a formal series. Let $\varphi$, normalised by
$\varphi(0)=0$ and $\varphi'(0)=1$, solve $\varphi(st) = R(\varphi(t))$. If
$\varphi$ is differentially algebraic over $\mathbb C(t)$, then
\begin{enumerate}[label=\textup{(\roman*)}]
\item the fixed point $0$ is \emph{repelling}, $|s|>1$; and
\item $\varphi$ is a Möbius transformation of one of
\begin{equation*}
\exp(\alpha t^{\varrho}),
\qquad
\cos(\alpha t^{\varrho}+\beta),
\qquad
\wp(\alpha t^{\varrho}+\beta),\ \wp^{2},\ \wp^{3},\ \wp'(\alpha t^{\varrho}+\beta),
\end{equation*}
where $\wp$ is a Weierstrass function, $\varrho\in\mathbb Q$ is such that the
function concerned is meromorphic on $\mathbb C$, $\alpha\neq0$, and $\beta$ is a
fraction of a period. Correspondingly $R$ is Möbius-conjugate to $z^{\pm d}$, to
$\pm T_d$ (Chebyshev), or to a Lattès map.
\end{enumerate}
\end{theorem}

\begin{remark}[Lineage and independent corroboration]
\label{a:rem:lineage}
For $|s|>1$ --- the Poincaré-function case --- the classification goes back to Ritt
in 1926 \cite{Ritt1926}; Becker and Bergweiler completed the picture for the three
classical functional equations at once, and Bergweiler's later account with
Aschenbrenner \cite{AschenbrennerBergweiler2015} assigns the exceptional list
explicitly to the repelling case. Di Vizio and Fernandes
\cite{DiVizioFernandes2021} have given an independent difference-Galois proof of
Ritt's theorem; notably, their Theorems~1.2 and~4.2 assume only that $s$ is neither
zero nor a root of unity --- no condition $|s|>1$ --- and show that any
differentially algebraic solution of $y(R(x)) = s\,y(x)$, which is our
$\psi$-convention with no inversion needed, must already satisfy a first-order
equation of Ritt type $(y')^{\varrho} = B(x)y^{\,j}$.
\end{remark}

\subsection{Hypertranscendence throughout the subcritical window}
\label{a:sec:htmain}

\begin{theorem}[No exceptional parameters in $(0,1)$]
\label{a:thm:ht}
For every $r\in(0,1)$, the Koenigs function $\psi_r$ of $f_r(z)=rz(1-z)$ at $0$ is
hypertranscendental over $\mathbb C(z)$: it satisfies no algebraic differential
equation. In particular $\psi_r$ is not an elementary function.
\end{theorem}

\begin{proof}
Fix $r\in(0,1)$. The map $f_r$ is rational of degree two, fixes $0$, and its
multiplier $s = f_r'(0) = r$ satisfies $0<s<1$; the hypotheses of \cref{a:thm:BB}
hold, with $0$ attracting. Suppose for contradiction that $\psi_r$ were
differentially algebraic. Its local inverse $\varphi_r = \psi_r^{-1}$ solves the
parametrisation form \cref{a:eq:schroderparam} with the stated normalisation, and by
\cref{a:lem:invda} it would be differentially algebraic too. \Cref{a:thm:BB}(i)
would then force $|r|>1$, contradicting $r\in(0,1)$. Hence $\psi_r$ is
hypertranscendental, and non-elementarity follows from \cref{a:lem:elemda}.
\end{proof}

\begin{remark}[The exceptional set inside the window is empty, not merely avoided]
\label{a:rem:empty}
The affine conjugacy $c(r) = (2r-r^2)/4$ of \cref{a:lem:affine} sends $r\in(0,1)$ to
$c\in(0,\tfrac14)$, which avoids the exceptional quadratic parameters $c=0$
(reached at $r=2$) and $c=-2$ (reached at $r=4$, and also at $r=-2$). The
exceptional Schröder functions exist only over \emph{repelling} fixed points, so the attracting window
contains no candidates whatsoever, and no enumeration of exceptional parameters is
required. Consistently, the classical elementary cases of \cref{a:app:integrable}
have multipliers $2$ and $4$ at the linearised fixed point --- both repelling: they
are Poincaré-type linearisations, and neither can be pressed into service for
$r\in(0,1)$.
\end{remark}

\begin{figure}[ht]
\centering
\includegraphics[width=0.66\textwidth]{figure1_parameter_conjugacy.pdf}
\caption{The affine parameter correspondence $c(r) = (2r-r^2)/4$ between the
logistic family $f_r(x)=rx(1-x)$ and the quadratic family $P_c(z)=z^2+c$. The
shaded segment is $r\in(0,1)$, the subcritical branching window, which maps to
$c\in(0,\tfrac14)$; the two integrable parameters, $r=2$ ($c=0$, monomial) and
$r=4$ ($c=-2$, Chebyshev), lie outside it.}
\label{a:fig:conjugacy}
\end{figure}

\subsection{Where the window sits in the quadratic family}
\label{a:sec:mandelbrot}

The main cardioid of the Mandelbrot set
--- the set of $c$ for which $z^2+c$ has an attracting fixed point --- is
parametrised by that point's multiplier $\lambda$ through
$c = \lambda/2-\lambda^2/4$, which is exactly the correspondence of
\cref{a:lem:affine} with $\lambda=r$. The branching window $r\in(0,1)$ is therefore
the real diameter of the main cardioid, traversed from the centre $c=0$ out to the
parabolic cusp $c=\tfrac14$, where $\lambda=1$ and the process turns critical.
\Cref{a:fig:mandelbrot} places the segment in context.

Throughout the window the Julia set is a quasicircle bounding the basin of the
attracting fixed point, and it is a genuine fractal for every $c$ in the interval.
The Koenigs function is the conformal coordinate that linearises the dynamics
inside that basin, so its natural domain is bounded by exactly this curve. A
function whose domain of definition has no smooth description is not a promising
candidate for a closed form, and $A(p)$, being one of its values, inherits the
difficulty. The contrast with the integrable cases is sharp: at $c=0$ the Julia set
degenerates to the unit circle, and it is precisely there that the linearisation
becomes elementary, \cref{a:eq:psi2}. The picture illustrates \cref{a:thm:ht}; it
does not prove it.

\begin{figure}[ht]
\centering
\includegraphics[width=0.92\textwidth]{figure4_mandelbrot_context.png}
\caption{The Mandelbrot set in the complex $c$-plane, with the interval
$c\in(0,\tfrac14)$ highlighted. Under the correspondence $c=(2r-r^2)/4$ of
\cref{a:lem:affine} this is the image of the subcritical branching window
$r=2p\in(0,1)$, and it forms the real diameter of the main cardioid, running from
the centre $c=0$ to the parabolic cusp $c=\tfrac14$, which is criticality. The two
integrable parameters are marked: $c=0$ (monomial, reached in logistic coordinates
at $r=2$) and $c=-2$ (Chebyshev, $r=4$), the latter at the far tip of the real
slice. Note that $c(r) = c(2-r)$, so each $c\in(0,\tfrac14)$ is the image of two
logistic parameters, $r$ and $2-r$; these correspond to the two fixed points of
$z^2+c$, with multipliers $r$ and $2-r$, and it is the attracting one, $r<1$, that
governs $A(p)$.}
\label{a:fig:mandelbrot}
\end{figure}
